
\begin{secao}{Atitude, bixes!}
\quadrinhos{7}

Na USP, os alunos têm a liberdade e apoio de se organizarem
para montar grupos de debates, ciclos de palestras, grupos
de desenvolvimento e até mesmo grupos para jogarem alguma coisa (como
RPG, Magic, Yu-Gi-Oh ou algum esporte).
Portanto, caso vocês tenham algum projeto em mente, não hesitem
em se organizarem com seus amigos e se informarem sobre como avançar com essa
ideia. Lembre-se de que seus veteranes estão aí para te aconselhar e tirar
eventuais dúvidas.

É possível também juntar-se com alguns amigos e formar grupos de
estudo, seja para alguma matéria com a qual vocês tenham dificuldade, para
discutir aquele EP/lista de exercícios que ninguém está conseguindo
fazer ou simplesmente estudar algum tópico de interesse mútuo.

A seguir, temos (majoritariamente) alguns dos exemplos dos grupos que foram 
direta ou indiretamente criados por alunos do nosso Instituto!

% RDs --------------------------------------------------------------------------
\begin{secao}{O que é RD?}

RD é o Representante Discente. É um forte elo de ligação entre professores e alunos. 
RD é um aluno que representa os nossos interesses frente aos diversos conselhos existentes.
O RD ajuda, junto com os conselhos, a decidir coisas como autorização para festas, 
mudanças no currículo, aumento de vagas na FUVEST, mudança no corpo docente (às 
vezes lutamos para tirar algum professor), enfim, coisas desse tipo e muitas mais. 

Acho que você já percebeu o quanto é importante ter um aluno em cada um desses conselhos. 
Infelizmente, não costumamos preencher todas as vagas que nos é de direito. Isso se 
deve ao desinteresse de alguns ou  falta de tempo da maioria de seus VETERANOS. 

É, bixo, qué você quem tem mais tempo para fazer as coisas funcionarem aqui, já que 
ainda não sabe o que é Rec, DP, Trabalho, Estágio etc. Portanto, se você quer fazer 
alguma coisa pelo lugar onde você estuda, está aí uma dica. Para você ser RD é 
necessário se candidatar. O mandato é de um ano.\footnote{Até 2008 ou por aí, as eleições de RD eram no primeiro semestre do ano. Agora, são geralmente no final, ou seja, como bixo você não poderá atuar como RD, mas candidate-se no final do ano!}. Observação: ser um RD é também uma boa maneira de saber como pensam os seus professores e como as coisas funcionam aqui.

Em 2012, excepcionalmente, as eleições serão feitas no comecinho do ano (13, 14 e 15 de março). De acordo com o edital (nos murais, emails, fique atento!), você, bixo, por mais que não possa se canidatar a nenhum cargo, pode exercer seu direito de voto! Procure conversar com seus VETERANOS para saber melhor como funcionam essas coisas. Por enquanto, vai aí um breve resumo do que mais ou menos acontece em cada um dos órgãos nos quais temos direito a representate(s).
  
No IME, temos 26 cargos de RD, sendo 10 necessariamente de
pós-graduação e 14 necessariamente de graduação (Os dois cargos
restantes são livres). Todos os cargos tem direito a um suplente.
 
Existem diferentes níveis de hierarquia na administração.
 
{\bf As CoCs,
Comissões Coordenadoras de Curso (Lic, Pura, Estatística, Aplicada e
Computação)} são as mais próximas dos alunos. Temos um cargo de aluno em cada comissão. São comissões
pequenas, que tratam
dos problemas internos de cada curso: mudança de currículo,
requerimentos, optativas. Subordinada à CG e ao conselho do relativo
departamento. Analogamente, temos um cargo em cada Comissão
Cordenadora de Programa (de Pós).
 
{\bf Os Conselhos de Departamento (MAT, MAE, MAC e MAP)} tem uma dinâmica
um pouco diferente das CoCs, são mais formais. Cada conselho
se reúne (quase) mensalmente e são formados (em geral) por mais pessoas,
sendo
que existem regras sobre participação dos diferentes níveis
hierárquicos de professores (Titular, Associado, Doutor e Assistente). Nesses
conselhos, além de aprovar algumas das decisões das Comissões
Coordenadoras de Curso e de Programa (pós) e distribuição de
carga didática, são discutidos re-oferecimento de curso, revisão de prova,
supervisão das atividades dos docentes, afastamentos (temporários ou não),
contratação de professores e muitas outras coisas.

Os Conselhos de Departamento são subordinados à Congregação e ao CTA.
 
{\bf A Comissão de Graduação (CG)}, basicamente, avalia requerimentos,
mudança/criação de cursos e jubilamentos.
Analogamente, existe a Comissão de Pós-Graduação (CPG). Ambas são
subordinadas à Congregação.
 
{\bf Comissão de Espaço Físico (COESF)} é um orgão consultivo do CTA, formado
por representantes de diversos "ocupadores de espaço": Biblioteca, Centro
de Software Livre, Matemateca. Também tem representantes de cada
departamento. É presidida pelo vice-diretor. O RD daqui é o mesmo do CTA.
 
{\bf A Comissão de Cultura e Extensão (CCEx)} quase nunca tem reunião. Cuida das
atividades de extensão: Matemateca, CAEM, etc...
 
Os dois conselhos mais importantes são o CTA e a Congregação, ambos
presididos pelo Diretor.
 
{\bf O Conselho Técnico e Administrativo (CTA)} cuida de todas questões não
acadêmicas: Orçamento, reformas, avaliação dos funcionários, xerox,
lanchonete. É formado pelos 4 chefes de departamento, diretor, vice diretor, um
representante dos funcionários e um RD.
 
{\bf A Congregação} é o órgão máximo do Instituto. Com muitos professores, a
maioria titular. São dois RDs de graduação e um de Pós. Basicamente, neste
órgão, são rediscutidas e aprovadas (ou não) muitas das decisões dos órgãos
subordinados. Os membros da Congregação tem voto na eleição para Reitor e
Vice-Reitor. 
 
Bom, bixo, caso você não tenha lido o começo desse texto, não é difícil perceber que é muito importante ter um aluno em cada um desses conselhos. Pergunte, participe, vote. Saiba do que anda acontecendo! 

%\quadrinhos3

\pagebreak
\end{secao}


% Rede Linux -------------------------------------------------------------------
\begin{subsecao}{Rede Linux}

\figurapequenainline{rede_linux}

\begin{subsubsecao}{Introdução}

A Rede Linux é uma rede de computadores, administrada por alunos do IME e
que fornece diversos serviços para os VETERANOS e até mesmo para vocês, bixos.
Ela disponibiliza:

%FIXME
\vspace{-1em}

\begin{itemize}
\item 2 salas de computadores (no bloco A) com todo\footnote{ Se um programa
estiver faltando, mande um email pra admin@linux.ime.usp.br pedindo-o.} tipo de
programa necessário para suas atividades acadêmicas (com pelo menos uma que fica
aberta 24 horas por dia, 7 dias por semana\footnote{ Mas talvez vocês não
consigam entrar no bloco A depois da meia noite, que é quando a portaria
fecha.});
\item Uma página na internet para cada aluno;
\item Um \textit{e-mail} para cada aluno;
\item Espaço para você guardar seus arquivos;
\item Acesso remoto via ssh (linux.ime.usp.br);
\item Impressoras;
\item Admins dispostos e capazes, para o caso de algum usuário ter alguma boa
ideia para adicionar a esta lista;
\end{itemize}
\end{subsubsecao}

\begin{subsubsecao}{O Linux}

A rede utiliza em todos os seus computadores um sistema operacional chamado
Linux. Esse é um sistema desenvolvido de forma colaborativa pelos usuários
e empresas interessados nele (se quiser saber mais a respeito, pesquise
por ``software livre''!).

O Linux não é um sistema mais difícil de usar que o Windows. É apenas
diferente em alguns aspectos. Além de tudo, existem cursos de Linux que são
organizados pelos alunos do IME. Os admins costumam promover esses cursos.
Fiquem atentos aos emails!

Não se deixem intimidar pelo sistema. Se vocês se derem ao trabalho de
aprender a utilizá-lo bem, verão que ele é bastante flexível, e até mesmo
interessante (tanto quanto um sistema operacional pode ser =P).

\end{subsubsecao}

\begin{subsubsecao}{Os admins}

Os admins são alunos do bacharelado em ciência da computação (vulgo BCC) que
são responsáveis por administrar a rede. Entre outras coisas, isso quer dizer
manter os computadores funcionando, ajudar os alunos a usar a rede (com
cursos\footnote{ Fiquem atentos aos emails!!!} e resolvendo dúvidas nos
horários de plantão\footnote{ Na página da rede, estão os horários de todos os
admins.}) e também implementar coisas novas na rede (aceitamos sugestões!)

Os admins são escolhidos por um treinamento que acontece de dois em dois anos,
em todo ano par. Mais informações serão divulgadas quando este estiver próximo
a ocorrer.

\end{subsubsecao}
\begin{subsubsecao}{Como criar uma conta?}

Basta passar na Admin, na sala 125 do bloco A (como vocês são bixos: bloco A é o da
biblioteca, bloco B aquele que tem muitas salas de aula e que vocês vão passar boa
parte da vida de vocês). Contatos:

%FIXME
\vspace{-1em}

\begin{description}
\item [e-mail:] admin@linux.ime.usp.br
\item [Página:] \url{www.linux.ime.usp.br}
\item [Sala:] 125, bloco A
\end{description}

%FIXME
\vspace{-.5em}

\end{subsubsecao}

\end{subsecao}


% FLUSP ------------------------------------------------------------------------
\begin{subsecao}{FLUSP}

\figurapequenainline{flusp}

Você já ouviu falar sobre Linux? E GCC? Talvez Gimp, TensorFlow, VLC ou Firefox?
Todos estes projetos são chamados de FLOSS: Free Libre Open Source Software, e
têm em comum a liberdade de código e conhecimento. Isto quer dizer que qualquer
um pode ler, usar e contribuir com código ou arte a estes projetos!

O FLUSP: FLOSS@USP é um grupo de extensão com o objetivo de reunir alunos de
graduação e pós-graduação interessados em contribuir para projetos FLOSS.
Atualmente temos contribuidores no Kernel Linux, no compilador GCC e muitos
outros projetos. Em 2018, o FLUSP foi responsável por 20\% das contribuições
para o Kernel Linux no subsistema Industrial Input/Output. Devido a nossas
contribuições para os drivers da Analog Devices, a empresa doou duas placas
de testes ao grupo.

Assim como na comunidade FLOSS, nós encorajamos a liberdade no FLUSP. Seja para
contribuir com um projeto existente ou compartilhar um projeto pessoal para
outros contribuirem com, aqui você encontra espaço!

\begin{description}
  \item[Facebook:] \url{facebook.com/flusp}
  \item[Telegram:] \url{http://tiny.cc/flusp}
  \item[IRC:] Servidor \texttt{irc.freenode.net}, canal \texttt{\#ccsl-usp}
  \item[Site:] \url{https://flusp.ime.usp.br}
  \item[Lista de email:] \texttt{flusp@googlegroups.com}
  \item[GitLab:] \url{https://gitlab.com/flusp}
\end{description}

\end{subsecao}


% IME Júnior -------------------------------------------------------------------
\begin{subsecao}{IMEjr: A Nossa Empresa}

\figurapequenainline{imejr_logo}

Em meados de 1991, surgia a Empresa Junior de Informática, Matemática e
Estatística do IME (IMEjr). Uma Empresa Junior é uma Associação Sem Fins
Lucrativos administrada por estudantes de graduação (que é o que você é agora)
e tem o objetivo de complementar a formação do aluno em termos da integração
entre teoria e prática, além de incentivar o empreendedorismo entre os alunos
do Instituto

Entre as atividades da IMEjr estão o desenvolvimento de projetos em todas as
áreas do IME e a organização de palestras, cursos e workshops. Logo, estamos
abertos tanto a alunos interessados em aprender a administrar uma empresa
quanto a desenvolver atividades e projetos.

Uma diferença entre nós e uma empresa comum é que nossos integrantes têm muito
mais liberdade de trabalhar e participam de projetos que auxiliam a aprofundar
mais sua formação, do que a maioria dos estágios por aí.

Contamos com vocês neste ano, e já garantimos que existem atividades prontas.
Esperamos seu contato! LEMBRE: nem sempre de aulas e livros é feito um
estudante com boa formação. Por isso, anote em sua agenda, celular, bloquinho
de nota, ...:

%FIXME
\vspace{-1em}

\begin{description}
\item [Sala:] 258, bloco A
\item[E-mail:] \url{imejr@ime.usp.br}
\item[Website:] \url{www.ime.usp.br/imejr}
\item[Facebook:] \url{fb.com/IMEJuniorUSP}
\end{description}

%FIXME
\vspace{-.5em}

\end{subsecao}

\pagebreak
% MaratonUSP -------------------------------------------------------------------
\begin{subsecao}{MaratonUSP}

\figurapequenainlineapertada{maratonusp}

O MaratonUSP é um grupo de estudos de programação competitiva. O principal evento é a Maratona de Programação, uma competição anual organizada pela Sociedade Brasileira de Computação (SBC). A competição é restrita a universitários que se organizam em times de três pessoas para resolver os mais variados problemas em forma de código. Apesar da coletividade, há uma dificuldade a mais: cada equipe dispõe de apenas um computador.

A Maratona é dividida em várias fases. Primeiro, há uma seleção interna dos times que representarão a USP na competição. Uma das vagas é reservada para um time de bixos! Depois, há uma fase regional com outras universidades da capital paulista, onde a USP tem um histórico de vitórias absoluto.

As coisas são consideravelmente mais difíceis na fase nacional, conhecida como Final Brasileira. As melhores universidades do Brasil se reúnem em algum canto do país para disputar as vagas brasileiras para a fase mundial. Exemplos recentes são Belo Horizonte, Foz do Iguaçu e Salvador. A hospedagem, comida e transporte são todos cobertos pela organização e pela universidade. Viagem de graça! Os melhores times brasileiros vão para a Final Mundial, sempre num país diferente. Nós estivemos recentemente no Marrocos, na Tailândia e na China. Em hotéis 5 estrelas! De graça!

Grandes empresas têm muito interesse nos alunos que participam da Maratona. É muito comum encontrar ex-maratonistas em empresas como Google, Facebook, etc. O pensamento abstrato e o trabalho em equipe desenvolvidos na Maratona são características muito bem avaliadas pelos recrutadores.

O ano começa com o bixeCamp, uma série de aulas focadas em ensinar o básico de programação e algoritmos. As aulas são no início dos treinos, que ocorrem no CEC (do lado da Seção de Alunos) à partir das 14h às sextas-feiras. O treino acaba às 19h, mas não é necessário ficar até o final. Alunos já experientes em programação são encorajados a comparecer também. Te esperamos lá!

\begin{description}
\item [Facebook:] \url{facebook.com/MaratonUSP}
\item[Site:] \url{www.ime.usp.br/~maratona}
\end{description}

\end{subsecao}

\pagebreak
% USPCodeLab -------------------------------------------------------------------
\begin{subsecao}{USPCodeLab}

\figurapequenainline{uspcodelab}

O USPCodeLab é um grupo de extensão que tem por objetivo criar um espaço
colaborativo para o desenvolvimento de tecnologia na USP.

O grupo foi fundado por estudantes do BCC em 2015 com o nome IME Workshop, mas
atualmente conta com integrantes de vários cursos da USP. Nosso foco é aprender
na prática ferramentas e técnicas de desenvolvimento de software que permitam
solucionar problemas do mundo real.

Durante o semestre, organizamos um grupo de estudos com duas turmas (iniciante
e avançada) que se reúnem semanalmente para estudar programação web/mobile.
Utilizamos os recursos mais avançados, criados por grandes empresas como Google
e Facebook e pela comunidade de software livre, para criar projetos legais
propostos pelos próprios participantes! Alguns exemplos são o nosso \textit{bot}
do bandejão (@uspbandexbot, no Telegram) e o sistema de reserva de armários do
CAMat.

O USPCodeLab também organiza o HackathonUSP, a maior competição de programação
da universidade. Durante o evento, os participantes se reúnem em times e são
desafiados a fazer, em apenas 24h, um protótipo de software ou hardware
relacionado ao tema da competição. Diversão garantida para quem gosta de comida,
café e muita programação!

Curtam nossa página do Facebook e entrem no nosso grupo do Telegram para saber
datas e horários das nossas reuniões. Participem do USPCodeLab!

\begin{description}
\item[Telegram:] \url{bit.ly/uspcodelab_tg}
\item[Reddit:] \url{reddit.com/r/uspcodelab}
\end{description}

\end{subsecao}


% USPGameDev: Pesquisa e Desenvolvimentos de Jogos na USP ----------------------
\begin{subsecao}{USPGameDev: Pesquisa e Desenvolvimentos de Jogos na USP}

Valve. Blizzard. Rockstar. Nintendo. USPGameDev. O que esses nomes têm em comum?
São nomes de grupos de desenvolvedores de jogos. E um deles tem sua sede na USP.

Constituído primariamente de alunos da USP de diversas áreas (na prática não) o 
USPGameDev (UGD) foi criado em 2009, já tendo lançado cinco jogos e 
publicado seu próprio \textit{kit} de desenvolvimento para jogos 2D e 3D que 
oferece suporte para diversas plataformas, tais como Windows, Linux, Mac e Android. 

Mas não é nada tão complicado assim. Mesmo tendo acabado de entrar na faculdade,
você também pode criar seu próprio jogo com a ajuda do USPGameDev. Tanto que nossos
últimos jogos lançados foram produtos de grupos de bixos com tempo livre demais em
suas vidas. Simplesmente apareça em uma de nossas reuniões e participe. Nenhum
conhecimento é necessário, até porque um dos principais objetivos do grupo é
aprendermos!

Além disso, o UGD também oferece cursos e \textit{workshops} para a comunidade USP
sobre diversos assuntos envolvendo desenvolvimento de jogos. Fique de olho!

Acesse nosso muito bem desenvolvido \textit{site}: 
\url{http://uspgamedev.org}

Para saber mais sobre os horários e datas das reuniões e como participar, acesse: 
\url{http://uspgamedev.org/contato}

\end{subsecao}



% IMEsec -----------------------------------------------------------------------
\begin{subsecao}{IMEsec}

\figurapequenainline{imesec}

O IMEsec é um grupo de extensão focado em aprender, estudar e se divertir com a
segurança da informação. Sem que a maioria das pessoas se dê conta, este nicho
está presente no cotidiano. Mandar uma mensagem no WhatsApp, navegar pela web,
entrar no Facebook: exemplos de ações simples do dia-a-dia que necessitam ser
feitas de maneira segura, visando à privacidade do usuário.

Nosso grupo, formado majoritariamente por alunos do IME-USP no início de 2017,
busca entender melhor este universo e expandí-lo no ambiente universitário. O
foco tem sido amplo; desde resolução de desafios on-line (que são muitos
divertidos), participação em competições (sim, competições — \textit{capture the
flag} — muito, muito legais) até desenvolvimento de projetos que possam
beneficiar a população. Todos são bem-vindos; basta ter interesse pelo assunto.

As experiências decorrentes de nossas atividades ajudaram e ajudam nos estudos,
além de agregarem valor à graduação. Por meio desses, mergulhamos não só em
computação mas também em matemática, estatística e até outras áreas bem
inusitadas.

Fiquem de olho na nossa página do Facebook e do Telegram para saber mais sobre
as reuniões semanais, além de fatos interessantes (ou simplesmente engraçados)
sobre o vasto mundo da segurança. Participem!

\begin{center}
  \Large
  \url{facebook.com/imesec}

  \url{Telegram: tiny.cc/imesec}

  \url{www.ime.usp.br/~imesec/}
\end{center}

\end{subsecao}


% Hardware Livre ---------------------------------------------------------------
\begin{subsecao}{Hardware Livre}

\figurapequenainline{hardwarelivre}

Nós, do Hardware Livre USP, desde 2013 nos encontramos semanalmente para
discutir sobre projetos de hardware. Estudamos e brincamos com diversos tipos de
hardware e software, desde Arduino até centrífugas de laboratório e impressoras
3D.

\textit{``Não sabemos programar e temos medo de tomar choque, podemos participar
do grupo?''}

Claro que podem! Provavelmente ao seu lado agora tem um microprocessador que
está pegando suas informações e processando de alguma maneira. Entender como
esses dispositivos funcionam é algo essencial para qualquer pessoa que queira
transformar o mundo em que vive construindo novos dispositivos para atender
suas necessidades.

\textit{``Ficamos sabendo que o grupo foi criado por alunos do BCC, não vamos
nos misturar com eles!''}

Não falem isso, um dos aspectos mais legais da universidade é a interação entre
diversas áreas do conhecimento e pessoas! No grupo já passaram pessoas de
diversos lugares, como da ECA, POLI, física, biologia… Até um pessoal do ITA
veio tirar umas dúvidas conosco. Então, se vocês não são do BCC, têm um motivo a
mais para participar.

Ao longo de todo o primeiro semestre, realizaremos uma série de workshops para
introduzir conceitos de programação e de hardware a quem nunca mexeu. Fiquem
ligados na nossa página no Facebook e no nosso site para saber mais! Mas vocês
também não precisam esperar até lá para nos encontrar — vejam em nosso site o
local semanal dos nossos encontros ou converse conosco no nosso grupo do
Telegram.

\begin{description}
  \item[Site:] \url{hardwarelivreusp.org}
  \item[Facebook:] \url{facebook.com/Hardwarelivreusp}
  \item[Telegram:] \url{tiny.cc/hardwarelivre}
\end{description}

\end{subsecao}

\pagebreak
% Tecs: Computação Social ------------------------------------------------------

\begin{subsecao}{Tecs}

\figurapequenainline{tecs}

O Tecs é um grupo de extensão que iniciou suas atividades no segundo semestre de
2017, e é focado no impacto social da computação e da tecnologia. 

Buscamos construir uma rede de estudantes interessados em melhorar a comunidade
local e a sociedade e criar oportunidades para que eles possam engajar-se em
projetos com esse fim; unir esforços para formar uma sociedade e profissionais
éticos e conscientes sobre o uso da tecnologia; promover a educação tecnológica
igualitária da população por meio de cursos, oficinas e ações promovidas pelo
grupo; e estimular alunos a usarem a tecnologia para solucionar problemas da
comunidade. 

Em termos gerais, pretendemos que os estudantes entendam como a tecnologia pode
ser utilizada para o bem coletivo, e utilizem esse conhecimento na prática, por
meio de colaborações com a comunidade local, os serviços públicos, as
organizações não-governamentais e as sem fins lucrativos. 

Se você tem interesse em promover o ensino de computação, em debater questões
éticas e sociais no contexto tecnológico ou em desenvolver aplicativos, sites ou
sistemas em parcerias com projetos sociais, entre em contato conosco e participe
do grupo!

\begin{description}
  \item[Site:] \url{www.ime.usp.br/~tecs}
  \item[Facebook:] \url{www.facebook.com/tecs.usp}
  \item[Telegram:] \url{https://t.me/tecsusp}
\end{description}

\end{subsecao}


%FIXME GAMBIARRA para não quebrar página num lugar zuado

% Diversime --------------------------------------------------------------------
\begin{subsecao}{DiversIME}

O grupo de Diversidade do Instituto de Matemática e Estatística tem como objetivo
 apoiar e unir a comunidade LGBT do IME. O grupo, fundado em 2012, possui hoje
 no facebook mais de 100 membros e nosso gráfico é do tipo $y = ax + b$ com $a < 0$,
 ou seja, crescemos constantemente.

A ideia principal é lutar contra o preconceito que ainda existe na sociedade,
 instruir as pessoas do instituto sobre o assunto e fazer amigos!

Esse ano, temos uma proposta de realizarmos reuniões mensais para estudar o
 assunto e nos apoiar mutuamente, seja em auxílio a estudantes com dificuldades
 na aceitação, com famílias que não aceitam ou até mesmo com a falta de
 namoradxs, o importante é participar.

E tudo isso também ocorre no nosso grupo secreto do facebook! É só pedir pra um
 veterano membro que ele te adiciona lá. Como o grupo é secreto, apenas membros
 podem te ver lá, garantindo sua segurança se você não se sentir confortável
 para abrir isso para a sociedade. Vale lembrar, é claro, que o grupo não
 abrange apenas LGBTs, mas também héteros amigos e interessados em apoiar
 a causa, participar de nossas discussões e causações. Tendo o coração e a mente
 abertos, são todos bem-vindos!

Sigam também nossa página no Facebook: {\tt fb.com/diversimeusp} !
Acredite, você não está sozinho, e a partir do DiversIME, você pode encontrar
 pessoas que passaram pelas mesmas dificuldades e maravilhas que você. Se você
 é imeano e compreende a importância da auto-aceitação e do respeito à
 diversidade, então você já é parte do DiversIME!

\end{subsecao}


% Existimos! -------------------------------------------------------------------
\begin{subsecao}{$\exists$xistimos!}

\figurapequenainline{existimos}

$\exists$xistimos! surgiu em  2014 com a proposta de criar um espaço de
confiança entre as alunas do IME, onde cada uma possa ter a liberdade e
segurança para discutir suas vivências e propostas.

O desconforto com a forma como as mulheres são tratadas no IME e nas ciências
exatas em geral fez com que nos juntássemos para discutir como as questões de
gênero se manifestam no instituto e quais são as formas de agir para evitar
situações desconfortáveis ou preconceituosas.

Desde então promovemos uma série de eventos e intervenções para que tal debate
atinja toda a comunidade imeana, além de realizarmos reuniões periódicas apenas
com mulheres para que possamos conversar e nos ajudar, em qualquer situação,
mas em especial naquelas onde possamos ser vítimas ou testemunhas de preconceito e
machismo. Todas vocês estão convidadas a participar das nossas reuniões, 
alunas de outros institutos, professoras, funcionárias e mulheres
em geral são muito bem vindas $<$3. 

Para falar conosco ou participar do grupo basta enviar um email
para 3xistimos@gmail.com ou existimos@google.groups.com ou ainda pelo facebook:
https://www.fb.com/3xistimos/

Para denúncias ou relatos anônimos acesse: http://bit.ly/existimos
\end{subsecao}


%FIXME GAMBIARRA para não quebrar página num lugar zuado


% Comissão de Acolhimento da Mulher! -------------------------------------------
\begin{subsecao}{Comissão de Acolhimento da Mulher - CAM}

A Comissão de Acolhimento da Mulher (CAM) foi criada no dia 9 de Novembro de 2016,
pela Portaria IME 1677, após a realização de várias reuniões abertas com amplos
debates durante praticamente um ano com o Coletivo Mulheres do IME, o qual 
apresentou a proposta ao Diretor do Instituto, tendo sido aprovada pela Congregação
da Instituição por unanimidade.

Somos uma Comissão assessora da diretoria do IME cuja a principal atribuição 
é dar acolhimento às vítimas de discriminação de gênero, de assédio moral e 
sexual e de violência contra a mulher, quando essas ocorrências envolverem 
pessoas da comunidade do IME ou tenham ocorrido em suas dependências.

Essas discriminações tornam a convivência mais difícil e provocam nas mulheres
a sensação de desrespeito e de não pertencimento. A criação de uma Comissão de 
Acolhimento da Mulher contribuirá para o combate institucional à violência contra
a mulher, à desigualdade de gênero e aos efeitos da cultura patriarcal na academia.

A Comissão é formada por duas professoras, duas funcionárias e duas alunas, 
eleitas por seus pares. Conheça a gestão atual: 
\pagebreak

\textbf{Professoras}: 
\begin{itemize}
  \item Lucia Junqueira (lucia@ime.usp.br)
  \item Renata Wassermann (renata@ime.usp.br)
\end{itemize}

\textbf{Funcionárias}: 
\begin{itemize}
  \item Gianne Uchôa (gianne@ime.usp.br)
  \item Luiza Ribeiro Camilo (luizinha@ime.usp.br)
\end{itemize}

\textbf{Alunas}: 
\begin{itemize}
  \item Catarina Braga (catarinabraga@ime.usp.br)
  \item Sophia Fiorotto (sophia@ime.usp.br)
\end{itemize}

As mulheres que procurarem a comissão poderão, se quiserem, indicar com qual ou 
quais de seus membros desejam conversar.

Para mais informações, acesse o site \url{https://www.ime.usp.br/~cam} ou veja
a página da CAM no facebook procurando por "Comissão de Acolhimento da Mulher - CAM" 
ou envie um e-mail para {\tt cam@ime.usp.br}. 


\end{subsecao}


%FIXME GAMBIARRA para não quebrar página num lugar zuado
\pagebreak

% CinIME -----------------------------------------------------------------------
\begin{subsecao}{CinIME}

Convidamos vocês a conhecer o CinIME!! Somos um projeto do CAMAT que promove o
lazer e a integração de IMEanos, entre docentes, discentes e funcionários por
meio de sessões de cinema em que vocês escolhem o filme!

O CinIME ocorre toda sexta-feira, às 16h, na sala B5. A sessão, o refrigerante
e a pipoca são grátis. Funciona assim: você entra na sala, se serve do que
quiser, senta, relaxa e se diverte assistindo ao filme da semana ao lado dos
seus amigos.

A escolha do filme é feita através de votação no grupo “CinIME USP”, no
facebook. Para participar, primeiro você acessa e curte a página do CinIME no
facebook (https://www.facebook.com/cinimeusp) para ter acesso às nossas
atualizações. Então, na aba “Sobre”, você clica no link de sugestões apontado
na descrição da página. Dentre todas as sugestões de filmes, cinco são
escolhidos pela equipe do CinIME para a votação, priorizando a diversidade de
gêneros cinematográficos. Depois disso, a votação é lançada no grupo “CinIME
USP”, no fim de semana, e se estende até a meia-noite da quinta-feira da semana
seguinte. O resultado é sempre anunciado ao fim da sessão da semana em curso,
que ocorre na sexta à tarde. O filme eleito é divulgado diariamente nos grupos
dos bixos e no grupo “IME-USP”, no facebook, e também por cartazes espalhados
pelo instituto. Uma vez por mês, a equipe do CinIME toma a liberdade de
escolher o filme da sessão sem votação.

Não se esqueçam de sugerir filmes, votar e comparecer ao CinIME!

Curtam nossa página no facebook e fiquem por dentro das novidades!

Parabéns e boa sorte nesta nova jornada, bixos!

\begin{description}
  \item[Facebook:] \url{https://www.facebook.com/CinIMEUSP}
\end{description}

\end{subsecao}


% Grupo A5 ---------------------------------------------------------------------
\begin{subsecao}{Grupo S4}

\begin{wrapfigure}{r}{0.25\textwidth}
    \vspace{-25pt}
    \begin{center}
      \includegraphics[width=0.25\textwidth]{img/logo S4 maior.png}
    \end{center}
    \vspace{-25pt}
  \end{wrapfigure}


Nós somos um grupo de estudantes de diferentes anos da graduação e pós-graduação 
do IME, e organizamos seminários feitos por alunes e para alunes em diversos
tópicos de matemática, muitos não vistos na graduação. Originado em 2015 na forma
de seminários de IC, o S4 tem hoje em dia o propósito de apresentar seminários 
semanais acessíveis e interessantes para estudantes que ainda não possuem muita 
bagagem matemática, inclusive de outros cursos da USP. 

Somos independentes de outras instituições do IME, e estamos sempre abertos a 
receber novos participantes, dos mais engajados aos mais casuais. É um espaço
para que es alunes tenham a oportunidade de apresentar algum tema que tenham
estudado, e começar na prática a treinar a habilidade de comunicação científica. 
Es alunes, independentes do ano, estão convidades a sugerir temas e dar ideias
para a organização de seminários. Caso desejem apresentar um seminário, nós
poderemos ajudar no seu preparo em um ambiente amistoso e receptivo.

Além disso, o S4 também é um ambiente de estudo e socialização, onde todes podem
compartilhar dúvidas, sugestões, conhecer novos bixes e veteranes, bater papo,
descontrair, etc. Estamos sempre dispostos a ajudar quaisquer alunes, especialmente
bixes, para que possamos ver um pouco mais do que a matemática tem a oferecer e
compartilhar as diferentes visões sobre ela de cada um de nós.

\begin{description}
  \item[Discord:] \url{discord.gg/XzXeCzc3rk}
  \item[E-mail:] \url{grupo-s4@riseup.net}
\end{description}

\end{subsecao}

% Olimpíadas de Matemática e Informática ---------------------------------------
% Maratona de Programação ------------------------------------------------------
\begin{subsecao}{Olimpíadas de Conhecimento}

\begin{itemize}

\item{\bf Matemática: }

Bom pessoal, se vocês entraram no IME, muito provavelmente já participaram
de alguma Olimpíada de Matemática no Ensino Fundamental e/ou Médio. A 
boa notícia é que vocês vão poder continuar participando se quiserem,
e quem nunca participou tem a oportunidade de começar agora.

Mas por quê participar? As Olimpíadas Universitárias de Matemática são uma
oportunidade de se divertir resolvendo problemas difíceis de Matemática e agregar
valor ao currículo ao mesmo tempo. Elas são parecidas com as Olimpíadas de
Ensino Médio, mas com conteúdo de Matemática da graduação (essencialmente 
Cálculo, Análise, Álgebra Linear, Álgebra, Combinatória e Teoria dos Números), 
mas com enfoque em problemas que exigem criatividade e técnicas mais inovadoras,
muitas das quais vocês provavelmente não verão durante toda a graduação.

De quais olimpíadas podemos participar? Como alunos de graduação, vocês podem
participar da Olimpíada Iberoamericana de Matemática Universitária (OIMU),
Olimpíada Brasileira de Matemática (OBM) e Olimpíada Internacional de 
Matemática (IMC).

Como fazemos para nos preparar? Os sites institucionais dessas olimpíadas
tem todo o material necessário para vocês que querem estudar e se preparar
para elas.

Como fazemos para participar? Inscrevam-se pelo site ou entrem em contato com
o professor Yoshiharu. Para o IMC aconselha-se ter ganhado medalha na OBM,
já que é necessário apoio financeiro do IME por ser uma olimpíada internacional.

Mas nós, bixos, temos chance? Como foi o desempenho de IMEanos nelas? Nós obtivemos
sucesso nestas olimpíadas. Ganhamos medalhas em todas as três competições e
o resultado mais recente foi uma medalha de bronze na IMC e ouro na OBM.

Links institucionais:

\url{http://www.cimat.mx/Eventos/oimu/}

\url{http://www.imc-math.org/}

\url{http://www.obm.org.br/opencms/}

Se tiverem alguma dúvida, não hesitem em perguntar a algum VETERANO sobre os
Olímpicos!

\item{\bf Informática: }

``Informática? Vocês mexem com Word, Excel e PowerPoint então?''

Responder essa pergunta já virou rotina para competidores da
Olimpíada Brasileira de Informática (OBI). Não, Informática 
não é Word. Oras, então o que é a OBI?

A OBI é uma competição de lógica e matemática. As provas são pequenos
problemas que você deve resolver com programas de computador.

Apesar de ser uma olimpíada que requeira um conhecimento mínimo de 
linguagem de computação, as provas em si geralmente não contém nada 
de teoria computacional mais avançada. 

Isso por que a OBI é voltada para alunos do Ensino Médio e recém
ingressantes. Quer dizer que vocês, bixos recém formados do Médio,
são nossa única esperança de ganhar medalhas e trazer glória ao IME!
Isso também quer dizer que essa é a única chance de vocês de fazer a OBI,
uma competição relativamente tranquila comparada à Maratona de 
Programação.

Para participar, basta falar com o Professor Carlinhos (\url{http://www.ime.usp.br/~cef/}), 
ou com ex-competidores, amáveis VETERANOS e lindas criaturas, que estarão disponíveis
durante a semana de recepção e em todos os treinos da Maratona de Programação também!

Para mais informações, acessem \url{http://olimpiada.ic.unicamp.br/}

\end{itemize}

\begin{subsubsecao}{Maratona de Programação}

À primeira vista, a Maratona de Programação pode soar um tanto
surreal. Nerds correndo pela USP ao mesmo tempo que resolvem
problemas de computação e matemática? Infelizmente esse não
é o caso.

A Maratona de Programação se resume à resolução de problemas.
Se você adora resolver desafios, quebrar a cabeça com novos
e excitantes problemas e acumular toneladas de dinheiro, esse
é o lugar perfeito para você!

A competição consiste de uma série de problemas que englobam
temas como programação dinâmica, grafos e estruturas de dados.
Times de três pessoas devem resolver a maior quantidade de
desafios em cinco horas de programação. E tudo isso com direito
a um lanche gratuito durante a prova.

Mas não temam, bixos. Não é só por que vocês acabaram de entrar que
a probabilidade de ganhar uma medalha seja nula. Inclusive, na primeira
fase da maratona, uma equipe de bixos tem vaga garantida para a
fase brasileira.

Além da fama, constantes pedidos por autógrafos e dinheiro de sobra,
a maratona também vai lhes trazer um conhecimento muito mais
adiantado em relação ao dos seus colegas de classe, e até oportunidades
de emprego em empresas de renome, como Google, Facebook e IBM.

Se vocês se interessaram pela maratona e querem saber os horários dos
treinos, como participar ou saber mais, acessem:

\url{http://www.ime.usp.br/~maratona}

Acessem também o \textit{site} da competição nacional:

\url{http://maratona.ime.usp.br/}

\end{subsubsecao}


\end{subsecao}


%FIXME GAMBIARRA para não quebrar página num lugar zuado
\pagebreak

% Fala Sério -------------------------------------------------------------------
\begin{subsecao}{Fala Sério}

\figurapequenainline{fala-serio}

A iniciativa Fala Sério é um grupo de extensão de cunho social da USP que atua em
escolas públicas de São Paulo proporcionando reflexões sobre futuro e felicidade.
O projeto se iniciou na Escola Politécnica da USP e devido ao seu grande sucesso 
se tornou um projeto de extensão da USP.

Para muitos alunos o processo de transição do Ensino Médio para a vida adulta é 
muito difícil e muitas vezes perturbador. Surgem muitas dúvidas, insegurança, 
cobranças; e é neste contexto que o Fala Sério atua.

Por meio de dois encontros levamos reflexões e espaços de conversa sobre 
felicidade e futuro; quais são nossos sonhos e como alcançá-los. Desse modo, 
oferecemos um apoio pessoal que muitas vezes os jovens não encontram em suas 
escolas e que faz uma grande diferença em suas vidas.

Eaí, se animou com a ideia? Quer compartilhar suas experiências e entrar nessa 
missão com a gente? Visite nossa página no facebook e vem nos conhecer:
\url{facebook.com/iniciativafalaserio}

\end{subsecao}


\end{secao}
